\subsection{An evaluated coherent-unit family at every positional level}

\begin{statement}{E.69}{Proved}{Kato's norm theorem in the actual complex tower and a product of global section fibers}
Keep $L_d=K(\beta_d)$ and the specified roots of E.64, where
$K=\mathbb Q(\zeta)$ and $\beta_d^{3^d}=17\zeta$.
Put
\[
 \mathcal U_d=\mathcal O_{L_d}[1/3]^\times
                         \otimes_{\mathbb Z}\mathbb Z_3,
 \qquad
 \mathcal U_\infty=\varprojlim_{d,N}\mathcal U_d,
 \qquad
 I_d=\operatorname{pr}_d(\mathcal U_\infty).
 \tag{E69.1}
\]
These are the usual finitely generated $3$-unit completions and their
compact inverse limit, with the actual adjacent field norms.
For every $d$, $I_d$ is a finite-index, torsion-free subgroup of
$\mathcal U_d$, free over $\mathbb Z_3$ of rank $3^d$.
The actual restricted norms give split exact sequences of compact
$\mathbb Z_3$-modules
\[
 0\longrightarrow K_d\longrightarrow I_{d+1}
      \xrightarrow{\ N\ }I_d\longrightarrow0,
 \qquad K_d\simeq\mathbb Z_3^{2\cdot3^d}.
 \tag{E69.2}
\]
Here the splittings are choices of lifts of an integral basis of
$I_d$; they are not asserted to extend to all of $\mathcal U_d$.

Let $\lambda=1-\zeta$. There is a unique finite integer $h\ge1$
such that
\[
                   I_0=\lambda^{3^h\mathbb Z_3}.
 \tag{E69.3}
\]
Its exact value, and the higher indices
$[\mathcal U_d:I_d]$, are retained as unevaluated finite norm-index
data. The assertions that they are finite and that $h\ge1$ are
proved below from the actual norm maps.
The first two finite projections are completely evaluated:
\[
 N_{L_1/K}\mathcal U_1=\lambda^{\mathbb Z_3},\qquad
 N_{L_2/K}\mathcal U_2=\lambda^{3\mathbb Z_3}.
 \tag{E69.4}
\]
For the second map an actual preimage of $\lambda^3$ is the
$3$-unit $c_1=\beta_1+3\zeta+1$, included into $L_2$.
The second image has index nine in
$\mathcal U_0=\mu_3\oplus\lambda^{\mathbb Z_3}$.
More generally define the finite projection exponent $a_n$ by
$N_{L_n/K}\mathcal U_n=\lambda^{3^{a_n}\mathbb Z_3}$ for $n\ge1$.
Then
\[
 a_1=0,\quad a_2=1,\quad
 a_n\le a_{n+1}\le a_n+1,
 \qquad a_n=h\ \hbox{for all sufficiently large }n.
 \tag{E69.5}
\]
Choose $s_d:I_d\to I_{d+1}$ with $Ns_d=\operatorname{id}$, and choose a basis
$k_{d,1},\ldots,k_{d,2\cdot3^d}$ of $K_d$.
The literal recursive maps
\[
 \begin{split}
 u_0&=\lambda^{3^h z},\\
 u_{d+1}&=s_d(u_d)
              \prod_{j=1}^{2\cdot3^d}k_{d,j}^{a_{d,j}},\\
 (z,(a_{d,j}))&\in
       \mathbb Z_3\times\prod_{d\ge0}\mathbb Z_3^{2\cdot3^d}
 \end{split}
 \tag{E69.6}
\]
give a topological group isomorphism onto the entire group
$\mathcal U_\infty$.
The inverse reads $z$ from $u_0$, and reads
$(a_{d,j})_j$ as the basis coordinates of
$u_{d+1}s_d(u_d)^{-1}$ in $K_d$.
In particular
\[
 \ker(\operatorname{pr}_0)=
       \prod_{d\ge0}\mathbb Z_3^{2\cdot3^d}
 \tag{E69.7}
\]
under these specified maps. More generally, every fiber of
$\operatorname{pr}_r:\mathcal U_\infty\to I_r$ is a torsor under
$\prod_{d\ge r}\mathbb Z_3^{2\cdot3^d}$.
The product in (E69.6) is an unrestricted compact product; its
finite-level coordinates are given recursively by the finite
products displayed there.

This gives actual global section classes through the following
closed embedding. Fix $A=3\kappa_3(17)$ and the origin
$g\mapsto(A_g,0)$ in E.46. For $u=(u_d)_d\in\mathcal U_\infty$
define the positional norm tuple
\[
 t_{d+1}=(1,u_d)\in
       \widehat{F_d^\times}_3\times\widehat{L_d^\times}_3,
 \qquad t_0=N_{K/\mathbb Q}u_0,
 \quad F_d=\mathbb Q(17^{1/3^d}).
 \tag{E69.8}
\]
All these tuples satisfy the total-norm cube condition defining
$J$ in E.46. Kummer--Shapiro therefore identifies them with precisely
the global $H$-section classes having pure component $1$ and complex
components in the $3$-unit groups at every level.
Their actual toric map (E46.7), in the coordinates (E69.6), is
\[
          (z,(a_{d,j}))\longmapsto
      \bigl(\kappa_3(17),\,
                   3^{h-1}z\,\kappa_3(3)\bigr).
 \tag{E69.9}
\]
Every nonempty fiber of this displayed family has the fully specified
product coordinates (E69.7).
In particular the fiber over $(\kappa_3(17),0)$ consists exactly
of the tuples with $u_0=1$ and is the full product in (E69.7).

The indices in (E69.3) have a further exact finite-level comparison:
\[
 I_d=\bigcap_{n\ge d}N_{L_n/L_d}\mathcal U_n,
 \qquad
 I_d=N_{L_n/L_d}\mathcal U_n
                 \quad\hbox{for every sufficiently large }n.
 \tag{E69.10}
\]
The proof supplies the actual maps and proves stabilization.
It does not compute the stabilization heights or identify a single
observed equality of consecutive images with permanent stabilization.
The remaining numerical indices and the sections with other
positional divisors remain separate data in the full norm system
of E.46. The maps from intrinsic metabelian and marked sections
to these $H$-sections remain the actual quotient maps of E.38;
(E69.8) supplies section classes in their target.
The manuscript anchor is \S51.16, scans172--178,
folios733v--736v, equations(371)--(379), and Part E(c),(d).
\end{statement}

\begin{proof}
\emph{The exact use of a coherent universal-norm theorem.}
Choose compatible roots of unity with $\zeta_3=\zeta$ and choose
the roots of seventeen so that
\[
 \beta_d=17^{1/3^d}\zeta_{3^{d+1}}.
 \quad
 \mathcal L=K(\mu_{3^\infty},17^{1/3^\infty})
       =L_d(\mu_{3^\infty},\beta_d^{1/3^\infty}).
 \tag{E69.11}
\]
Both equalities are literal equalities inside the chosen algebraic
closure. For example, dividing
$\beta_{d+j}$ by $\zeta_{3^{d+j+1}}$ recovers
$17^{1/3^{d+j}}$; multiplying these two factors recovers
$\beta_{d+j}$. Thus the two sets of generators in (E69.11) give
mutual field inclusions.

Kato defines universal norms as the image of the inverse limit of
completed $3$-units over all finite intermediate fields, with norm
transitions. His Theorem1.1(2) states that this image has finite
index in the base $3$-unit group when the base has exactly one prime
over three, for the radical--cyclotomic extensions in (E69.11)
\cite[pp.551--552, definition and Theorem1.1; Corollary5.2(2)]
{E_KatoUniversalNorms}.
The complete proof was checked through the cohomological norm
comparison in \S2.2--2.6, the abelian case Proposition3.1, the
Tor-module argument Proposition4.2, and \S5.4--5.5.
In particular the source theorem concerns compatible systems,
not just separate membership in finite norm images.

By E.64, $L_d$ has exactly one prime over three. Apply this theorem
with base $L_d$ and radical $\beta_d$.
There is the following actual restriction map of inverse systems:
\[
 \varprojlim_{L_d\subset E\subset\mathcal L\atop [E:L_d]<\infty}
       \bigl(\mathcal O_E[1/3]^\times\otimes\mathbb Z_3\bigr)
       \longrightarrow \mathcal U_\infty.
 \tag{E69.12}
\]
For $n\ge d$ its $n$th coordinate is the given coordinate at $L_n$;
for $n<d$ it is $N_{L_d/L_n}$ of the coordinate at $L_d$.
Norm transitivity proves compatibility at every adjacent arrow,
including the arrow across $d$. The projection of (E69.12) to
$\mathcal U_d$ is exactly Kato's universal-norm image.
Consequently this finite-index subgroup is contained in $I_d$.
The latter is closed, since it is the continuous image of the compact
group $\mathcal U_\infty$.
It follows that $I_d$ also has finite index in $\mathcal U_d$.
This proves the claimed finite-index assertion through a map of
the actual norm systems.

\emph{The finite ranks and the local torsion exclusion.}
The field $L_d$ has degree $2\cdot3^d$, has no real embeddings, and
has one prime over three. Dirichlet's unit theorem gives
\[
 \mathcal U_d\simeq\mu_3\oplus\mathbb Z_3^{3^d}.
 \tag{E69.13}
\]
For the roots of unity, E.64 proves the exact intersection with the
$3$-power cyclotomic fields: over any such field containing $K$,
$X^{3^d}-17\zeta$ remains Eisenstein at a prime over seventeen.
Hence $L_d$ is linearly disjoint from that field over $K$ and has
no additional $3$-power roots of unity. Its $3$-primary torsion is
precisely the displayed $\mu_3$.

Let $\mathfrak t_d$ be the unique prime over seventeen.
Reduction of $3$-units gives a continuous map
\[
 r_d:\mathcal U_d\longrightarrow
      (\mathbb F_{289}^{\times})^{\wedge_3}\simeq\mu_9.
 \tag{E69.14}
\]
The extension $L_{d+2}/L_d$ at these primes is totally ramified of
degree nine, with the same residue field. Its residue unit norm is
the ninth power. Thus
$r_d(N_{L_{d+2}/L_d}v)=r_{d+2}(v)^9=1$ for all
$v\in\mathcal U_{d+2}$.
Every $u_d\in I_d$ comes from that actual coordinate of a compatible
system, so
\[
                         I_d\subset\ker r_d.
 \tag{E69.15}
\]
Reduction is injective on $\mu_3$ because the residue characteristic
is seventeen. Equations (E69.13)--(E69.15) therefore prove that
$I_d$ has no torsion. A finite-index torsion-free submodule of the
finitely generated module (E69.13) is free of rank $3^d$.
This proves the precise assertion needed for (E69.2).

\emph{The actual base coordinate and the total-norm condition.}
The certified E.57 calculation gives the entire $3$-unit group at
the first complex level: its free basis is
$e_1,e_2,c_1=\beta_1+3\zeta+1$, its torsion basis is $\zeta$, and
their norms to $K$ are $1,1,\lambda,1$.
Indeed $(c_1)$ is the single prime over three, so adjoining $c_1$
to the ordinary unit basis gives all $3$-units.
It follows that
\[
       N_{L_1/K}\mathcal U_1=\lambda^{\mathbb Z_3},
       \qquad I_0\subset\lambda^{\mathbb Z_3}.
 \tag{E69.16}
\]
In $\mathbb F_{289}$ one has
$\lambda^2=-3\zeta$ and
$\lambda^{32}=(-3)^{16}\zeta^{16}=\zeta$.
The projection to the $3$-primary factor of this group of order
$288=32\cdot9$ is the power $64$; hence
$r_0(\lambda)=\zeta^2$.
Combining (E69.15) with (E69.16) shows that an exponent occurring
in $I_0$ is divisible by three.
The subgroup $I_0$ is closed and of finite index in
$\mu_3\oplus\lambda^{\mathbb Z_3}$, so it is exactly
$\lambda^{3^h\mathbb Z_3}$ for one finite $h\ge1$.
This proves (E69.3) with both the finite-index and lower-bound
assertions justified. It leaves the value of $h$ as the stated
finite arithmetic datum.

The same local argument already puts
$N_{L_2/K}\mathcal U_2$ in $\lambda^{3\mathbb Z_3}$.
For the opposite inclusion, the included element $c_1$ has
$N_{L_2/L_1}c_1=c_1^3$, and E.57 gives
$N_{L_1/K}c_1=\lambda$. Thus its total norm is $\lambda^3$,
proving the second equality of (E69.4) with an actual preimage.
The identity for its first norm can also be checked directly:
$N_{L_1/K}(\beta_1+1+3\zeta)
=(1+3\zeta)^3+17\zeta=1-\zeta$.
At every finite height the included $c_1$ supplies the nonzero
image $\lambda^{3^{n-1}}$; hence that closed image in
$\lambda^{\mathbb Z_3}$ has the unique exponent $a_n$.
For an element of $\mathcal U_n$ included into $\mathcal U_{n+1}$,
the adjacent norm is its cube. Therefore the next projection image
contains the cube of the current image and is contained in the
current image. This proves the inequalities in (E69.5).
Their eventual value is $h$ by (E69.10), proved next.

\emph{Surjective stable-image transitions and recursive coordinates.}
For $u_d\in I_d$, choose a compatible sequence that projects to it.
Its next coordinate lies in $I_{d+1}$ and norms to $u_d$.
Thus the actual map $N:I_{d+1}\to I_d$ is onto.
Choose any integral basis of $I_d$, lift its basis vectors to
$I_{d+1}$, and extend $\mathbb Z_3$-linearly. This gives the continuous
splitting $s_d$. Its kernel is free of rank
$3^{d+1}-3^d=2\cdot3^d$, proving (E69.2).

For the product on the right of (E69.6), the displayed recursion
uses a finite product at each level and has norm $u_d$ by construction.
It therefore gives an element of $\mathcal U_\infty$ for every
unrestricted set of product coordinates.
Conversely a compatible sequence has $u_0\in I_0$ and
$u_{d+1}s_d(u_d)^{-1}\in K_d$, so the inverse formulas determine
all its coordinates uniquely.
Each finite output coordinate depends continuously on finitely
many input coordinates, and the inverse basis maps are continuous
between finite free compact $\mathbb Z_3$-modules.
This proves the topological isomorphism, (E69.7), and the analogous
fiber formula at any level $r$.

To prove (E69.10), the inclusion from left to right is norm
compatibility. If an element lies in every norm image on the right,
then every finite set of the norm equations above level $d$ has a
solution: choose a preimage at its highest level and norm it to
the lower levels. These equations define closed subsets of the
product of the compact $3$-unit groups. The finite intersection
property gives a compatible infinite system.
Hence their intersection is $I_d$.
Every norm image contains the finite-index subgroup $I_d$; after
quotienting by $I_d$, these images form a descending chain of
subgroups in the finite group $\mathcal U_d/I_d$.
The chain stabilizes, proving the second assertion of (E69.10).
This argument uses compactness only for the finitely generated
$3$-unit completions in (E69.1).

\emph{The closed map to global sections and its toric coordinate.}
The natural map
$\mathcal U_d\to\widehat{L_d^\times}_3$ is a closed embedding.
For an elementary verification, a $3^r$th power in $L_d^\times$
which is a $3$-unit has a $3$-unit root, since all valuations away
from three of that root must vanish.
Thus the induced power filtration on the $3$-unit group is its
own multiplication filtration. The quotient is a subgroup of the
free divisor group away from three and is $3$-torsion-free;
the power quotient sequences and their inverse limits preserve
the asserted injection. Compactness of $\mathcal U_d$ and the
Hausdorff target make its image closed.

The positional algebra at level $d+1$ is $F_d\times L_d$, with
the root-cubing norms proved in E.46 and E.62.
Consequently (E69.8) is a continuous injection into the actual
positional norm limit. The norm of its base coordinate is
\[
           t_0=N_{K/\mathbb Q}(\lambda^{3^h z})
                   =3^{3^h z}.
 \tag{E69.17}
\]
Since $h\ge1$, this is a cube in
$\widehat{\mathbb Q^\times}_3$. It is exactly the condition
in E46.4 for the $J$-cocycle, with no additional lift chosen at
any positional level.
The E.46 Kummer--Shapiro map and its inverse identify this tuple
with an actual global section class and preserve the given
topology. The domain is compact and the target norm limit is
Hausdorff, so this is a closed embedding.
Conversely the specified pure and complex coordinates of any
section in this image give exactly a sequence in
$\mathcal U_\infty$; hence the image is the full family described
in the statement.
Finally E46.7 divides the two raw augmentation coordinates by
three. Applying Kummer to (E69.17) gives (E69.9).
The class $\kappa_3(3)$ has infinite order, as detected by its
valuation at three, so its zero fiber is exactly $z=0$, and the
other fibers follow from (E69.7).
\end{proof}
